\section{Úvod}

Než se budeme zabývat samostatnou \textbf{analytickou geometrií},
je třeba se napřed seznámit s tím, co to vlastně je.

Analytická geometrie je část geometrie, 
která se zabývá studiem geometrických útvarů pomocí 
algebraických metod. Tedy jednoduše řečeno, pomocí čísel a algebraických operací,
popisujeme a zkoumáme geometrické objekty jako jsou \textbf{body}, \textbf{přímky}, \textbf{kružnice}, \textbf{elipsy} a další.

To nám umožňuje řešit geometrické problémy pomocí algebraických rovnic a nerovnic,
což je často mnohem jednodušší než řešit je pomocí čistě geometrických metod.
Pro pochopení analytické geometrie je důležité se nejprve seznámit s \textbf{kartézskou soustavou souřadnic},
která nám umožňuje přiřadit každému bodu v rovině nebo prostoru jednoznačnou sadu čísel, které nazýváme \textbf{souřadnicemi}.